\section{Related work}
\label{sec:related}

Recall from \S\ref{sec:intro} that simulators are best compared by
\emph{altitude}: what one simulated run
represents. A run of \texttt{fleetsim} represents weeks of a whole fleet ---
many clusters, many tenants, thousands of jobs arriving, queueing, running,
being preempted and occasionally dying with their nodes. Most existing tools sit
at a different altitude, which is not a deficiency in them: a tool that models
one job's network traffic in microseconds is not a worse fleet simulator, it is
not a fleet simulator at all, and it answers questions \texttt{fleetsim} cannot.
Table~\ref{tab:related} summarises the landscape. One caveat: this is a survey of
what each tool's papers and documentation describe, not a benchmark --- we did
not install, run or measure any of these systems.

\begin{table}[!t]
\centering
\small
\caption{Related tools by altitude. Characterisations are taken from each
tool's published description, not from measurements we made.}
\label{tab:related}
\begin{tabular}{@{}p{2.6cm}p{2.2cm}p{4.0cm}p{4.3cm}@{}}
\toprule
\raggedright Family & \raggedright Altitude & \raggedright What it does well & \raggedright What a fleet question needs that it omits \tabularnewline
\midrule
\raggedright Blox \cite{blox2024}; the simulators inside Gavel, Pollux, Sia
  & \raggedright One cluster
  & \raggedright Fast, faithful replay of published deep-learning scheduling policies; Blox reports seven schedulers in 12--1157 lines each
  & \raggedright Capacity is a flat GPU count: no topology levels, no node failures or repair, no inference, no multi-cluster \tabularnewline
\addlinespace
\raggedright ASTRA-sim \cite{astrasim2023}, SimAI \cite{simai2025}
  & \raggedright One job $\times$ network
  & \raggedright Collective communication and per-iteration time at packet and flow detail
  & \raggedright There is no queue: the job is already placed, so there are no arrivals, no waiting and no preemption \tabularnewline
\addlinespace
\raggedright Vidur \cite{vidur2024}, LLMServingSim
  & \raggedright One inference service
  & \raggedright Request batching, KV-cache behaviour and tail latency for LLM serving
  & \raggedright Inference only, one service; training gangs, quota and failures are outside the model \tabularnewline
\addlinespace
\raggedright Batsim \cite{batsim2016}, the UBCCR Slurm simulator
  & \raggedright HPC batch
  & \raggedright Decades of batch-queue semantics, backfill \cite{easybackfill1995}, and a clean engine/policy separation
  & \raggedright No accelerator cost model: preemption is not priced in lost training work, and there is no checkpoint or gang-restart notion \tabularnewline
\addlinespace
\raggedright AIReSim \cite{aireesim}
  & \raggedright Reliability of one job
  & \raggedright Failure, repair and restart accounting for a single long training run
  & \raggedright No scheduler, no queue, no competing tenants \tabularnewline
\addlinespace
\raggedright kube-scheduler-simulator \cite{kubescheduler}
  & \raggedright Kubernetes placement
  & \raggedright Explains, decision by decision, why a pod landed where it did
  & \raggedright Wall-clock, no simulated time: no job completion time, no occupancy over a horizon \tabularnewline
\addlinespace
\raggedright Borg \cite{borg2015,borgtng2020}, MAST \cite{mast2024}, Singularity \cite{singularity2022}
  & \raggedright Fleet --- the existence proof
  & \raggedright Run real fleets at exactly this altitude and publish what the hard problems are
  & \raggedright They are production systems, not simulators; there is no public artefact in which to run an experiment \tabularnewline
\bottomrule
\end{tabular}
\end{table}

\paragraph{Single-cluster scheduler simulators.}
Blox \cite{blox2024} and the simulators built inside scheduling papers such as
Gavel, Pollux and Sia take a job trace and a policy, apply them to one cluster
of $N$ interchangeable GPUs, and report completion times. They are the closest
prior art, and \texttt{fleetsim} borrows from them directly: Blox's decomposition
of a scheduler into an \emph{ordering} policy (who runs next) and a
\emph{placement} policy (where it goes) is the API we adopted, and its convention
of measuring the steady middle of a run is our metrics window. What a flat GPU
count cannot express is a job failing to fit for structural reasons. Our
placement study makes the consequence concrete: on a 256-chip fleet at one seed,
changing only the placement rule left occupancy identical to the last digit
(0.8141 under three of four policies) while the mean queue wait of eight-node
gangs moved from 19{,}057\,s to 16{,}195\,s. A simulator whose capacity is an
integer cannot represent that mechanism; node failures are the second gap.

\paragraph{Per-job network simulators.}
ASTRA-sim \cite{astrasim2023} and SimAI \cite{simai2025} simulate one training
job's iteration in detail --- collectives, fabric topology, congestion. This is
precisely the layer \texttt{fleetsim} collapses into a single number: a
placement spanning a configured topology level multiplies the job's progress
rate by a configured multiplier, and there is no default multiplier because the
honest value is fleet-specific. The two altitudes compose rather than compete ---
a network-level study is how a lab would derive that multiplier. What these tools
do not have is a queue: one already-placed job has nothing to wait behind.

\paragraph{Inference-serving simulators.}
Vidur \cite{vidur2024} and LLMServingSim model a serving deployment: request
arrivals, batching, KV cache, latency percentiles. \texttt{fleetsim} models
inference at replica granularity --- a replica being a fixed block of chips held
for a service --- and stops there, because per-request burstiness sits below
fleet altitude.

\paragraph{HPC batch simulators.}
Batsim \cite{batsim2016} and the Slurm simulator carry decades of batch-queue
semantics, and \texttt{fleetsim} copies Batsim's central architectural choice:
the engine mutates state, policies emit intents, and the engine validates and
applies them, in-process rather than over a socket. EASY backfill
\cite{easybackfill1995} is a shipped \texttt{fleetsim} scheduler for the same
reason. The gap is not gangs as such --- MPI jobs are gangs --- but the
accelerator cost model around them: killing a training job destroys the work done
since its last checkpoint, a GPU job cannot be suspended in place because its
memory is held, and a restart costs minutes. Those costs are what make preemption
policy interesting, and what these models do not price.

\paragraph{Production fleet managers.}
Borg \cite{borg2015,borgtng2020}, MAST \cite{mast2024} and Singularity
\cite{singularity2022} operate at fleet altitude, and supply much of
\texttt{fleetsim}'s substance: Borg's priority bands are the tier semantics, and
published reliability figures \cite{metareliability2024,llama3} set the failure
defaults. Borg TNG's tail of per-job \emph{resource-hours} ($\alpha \approx
0.69$--$0.77$, the top 1\,\% of jobs accounting for 97--99\,\% of load) is a
target the traffic generator is meant to satisfy --- not a configured parameter,
and not yet checked by a shipped rung. What
none of them ships is something you can run, and the public traces
\cite{philly2019,helios2021,alibabapai2022,acme2024} record only what happened
under one policy, never what a different one would have done. That counterfactual
is the gap. One precision about the borrowing: \texttt{fleetsim}'s quota model is
inspired by MAST but is not a model \emph{of} it (Section~\ref{sec:limitations}).

The workload generator likewise stands in the workload-modelling tradition of
rigid parallel jobs \cite{feitelson2003}, with non-Poisson arrivals
\cite{paxsonfloyd1995} and lognormal duration bodies taken from the Helios and
Acme characterizations \cite{helios2021,acme2024} --- and the closed-form checks
of
Section~\ref{sec:validation} are ordinary queueing theory \cite{kingman1961}. A
cross-simulator check against Blox's published percentiles was planned and never
built.

\section{Limitations and future work}
\label{sec:limitations}

\paragraph{No fractional GPUs.}
The smallest unit \texttt{fleetsim} allocates is one whole chip. Alibaba PAI's
headline results \cite{alibabapai2022} --- roughly a 50\,\% saving from GPU
sharing, with a median instance requesting 0.042 of a GPU --- are therefore out
of reach, and no placement policy substitutes for the missing mechanism. It was
a target of the previous release that the release did not reach, and we record
it as still open rather than quietly re-dating it.

\paragraph{No per-job delay-cause attribution.}
\texttt{fleetsim} records how long each job waited, not \emph{why}. Philly's
split of queueing delay into a fair-share component and a fragmentation component
\cite{philly2019} --- the result that motivated this design, where fragmentation
accounts for 59--98\,\% of multi-GPU delay --- cannot be reproduced, and is a
second promise carried over from the release before last, still open. The metric
we did add, a count of
partially-occupied nodes, is a down payment, not a substitute, and it is
directional rather than proportional: an 8.9\,\% reduction in its
time-average accompanied a 25.7\,\% reduction in mean FIFO completion time,
because a mean over all instants dilutes exactly the moments when the queue head
needs a whole free node.

\paragraph{The performance model is analytical.}
A job's progress rate is the product of configured per-level multipliers. The
identity \emph{completion $=$ start $+$ overhead $+$ remaining work / rate} is
pinned exactly in integer microseconds, but the rate is an input, not a
derivation, and there is no default crossing penalty. Chip heterogeneity is worse
off: the throughput hook always returns 1.0, and scenarios using unpinned chip
types are rejected outright. CPU contention and co-location interference are not
modelled at all.

\paragraph{Failures are sampled from a rate, not from a physical model.}
One exponential draw at the fleet-wide aggregate rate, a thinning test, then a
uniform victim pick --- at a default mean time between failures of 42 node-days
with a fixed cause mix (60\,\% GPU/HBM, 10\,\% network, 13\,\% software, 17\,\%
other). Nothing in that is physical: no correlated failures (a rack losing power
takes one node at a time here), no thermal or age dependence, and unreliable
``lemon'' nodes are off by default. Spot preemption is deliberately harsher than
the products it abstracts --- zero notice, no banked progress --- so spot losses
are an upper bound, not an estimate.

\paragraph{One trace family, validated in depth.}
The quantitative validation is Helios \cite{helios2021}: four clusters, both
schedulers, the full September window. Philly \cite{philly2019} is exercised only
at converter level, and even there the full check, though written, was never run
against the real one-gigabyte artefact, so every Philly row reports
``unmeasured'' rather than a number; only structural assertions on a synthetic
2{,}000-row slice run in CI. Alibaba PAI is registered but never replayed. One
trace family is one lab's workload, one scheduler lineage, one hardware
generation.

\paragraph{The absolute numbers are order-sensitive.}
On Saturn, 35.5\,\% of the 105{,}969 jobs in the window share an exact submit
second, and 92.9\,\% share a 60-second scheduler round --- so the order in which
tied arrivals are served is a modelling choice, not a property of the trace.
Flipping the FIFO tie-break from ascending to descending job id moves Saturn's
mean FIFO completion time from 75{,}329\,s to 62{,}549\,s, a 17.0\,\% move, and
inverts which placement policy looks closer to published. Ascending id is
demonstrably the faithful order (ids are assigned at submission and are monotone
in submit time on three of the four clusters), so this is a sensitivity probe
under a deliberately wrong order, not an error bar. Its magnitude still matters:
17\,\% is wider than the $\pm 14$\,\% agreement we report on absolute completion
times, and a large fraction of the 25.7\,\% move the placement change itself
produced. No single cluster's absolute number can select a modelling choice,
which is why our discriminator aggregates four clusters --- and whether that
discriminator survives the same perturbation is \emph{unmeasured}: the probe ran
on one cluster under one scheduler (Saturn under FIFO, in both tie-break
orders), and extending it to all four clusters under both schedulers needs 24
more replays.

\paragraph{Single-metro scheduling only.}
There is one scheduling stage over one metro. Two-stage scheduling across metros
--- the MAST shape \cite{mast2024}, where a first stage decides which region a job
belongs in --- is schema-present and rejected as not implemented, as are TPU
optical-switch predicates, multi-gang jobs and autoscaling inference. The quota
model is also simpler than the systems that inspired it: \texttt{fleetsim}
commits quota at admission against queued demand and demotes irreversibly,
whereas MAST and comparable systems evaluate against running usage at scheduling
time and treat over-quota work as borrowed capacity that can be handed back. The
quota scenario shipped as \texttt{examples/06\_economics/} is not among the case
studies of \S\ref{sec:cases}; its demotion counts are partly an artefact of that
admission-time choice, which is why we do not quote them here.

\paragraph{Smaller ones, stated plainly.}
The backfill scheduler implements the conservative one-rule form of EASY, so
measured backfill rates undershoot published EASY reproductions. Byte-identical
output holds only within a platform, since the floating-point samplers differ in
the last bit across operating systems. And the performance envelope in the design
document --- six simulated months of a 100{,}000-GPU fleet --- is a target, not a
measurement: what has been measured is six months of a 2{,}048-chip fleet in
about nine minutes, and two days of a 524{,}288-chip fleet in about eighty
seconds, on a laptop.

\paragraph{Future work.}
Five items, in priority order. \emph{Fractional GPUs}, the single
change that unlocks a whole trace family's published results. \emph{Per-job
delay-cause attribution}, which would reproduce Philly's split and turn the
stranding metric into a per-job account of where a wait came from. \emph{More
trace families} --- the Philly artefact, then Alibaba PAI --- because one
validated family is the weakest part of Section~\ref{sec:validation}. \emph{TPU
torus geometry}: TPU pods wire their chips into a three-dimensional torus --- a
grid whose edges wrap around --- and a job is given a rectangular sub-block of
it, a \emph{slice}. Feasibility then depends on which free \emph{cubes}
(contiguous $4 \times 4 \times 4$ blocks) exist rather than on how many chips are
free, and the published optical-switch result removes contiguity constraints for
cube-multiple slices \cite{tpuv4isca2023}. And a
\emph{throughput matrix for heterogeneous chips}, replacing the hook that returns
1.0, so a fleet mixing chip generations can be modelled rather than approximated
as uniform.

\section{Conclusion}
\label{sec:conclusion}

We built an open-source discrete-event simulator that operates at fleet altitude:
a hierarchy of clusters, pods, racks and nodes; gang scheduling with atomic,
immutable allocations; priority preemption priced in lost work; node failures,
repair and maintenance drains; and a generative traffic model drawn from
published trace characterisations.

We then tried to break it. Against the four Helios clusters \cite{helios2021},
\texttt{fleetsim} reproduces the published ratio between FIFO and
shortest-job-first mean completion times with a mean absolute error of 3.4\,\%,
and lands within $\pm 14$\,\% of the published absolute completion times --- a
weaker, reported-not-asserted result the repository deliberately calls ``right
ballpark'' rather than reproduction. Both hold only under three modelling
choices we name every time we state the result, and inside tolerance bands wide
enough that a deliberately bad placement policy passes them too. The engine also
reproduces four textbook queueing results it could have failed --- Erlang-C,
Pollaczek--Khinchine, two-class preemptive-resume priority and
shortest-processing-time optimality --- plus a behavioural property test of EASY
backfill. And we have written down what it does not reproduce, and why.

If one sentence survives: matching a published table was not the useful part ---
the useful part was that matching it forced three modelling choices into the
open, and that occupancy, the metric everyone reports, was bit-for-bit identical
across placement policies whose queueing behaviour was not.
